\documentclass[12pt,a4paper]{ctexart}
\usepackage[top=2cm,bottom=2cm,left=2.5cm,right=2.5cm]{geometry}
\usepackage{listings}
\usepackage{xcolor}
\usepackage{hyperref}
\lstset{
  basicstyle=\ttfamily\small,
  breaklines=true,
  frame=single,
  numbers=left,
  numberstyle=\tiny,
  backgroundcolor=\color{gray!5},
  keywordstyle=\color{blue},
  commentstyle=\color{green!40!black},
  stringstyle=\color{red},
}
\title{LaTeX Test}
\date{\today}
\begin{document}
\maketitle
\tableofcontents
\newpage
\section{矩阵取数游戏}
\begin{tabular}{|l|l|}
\hline
队员 & 王梓豪 \\ \hline
日期 & 2026-07-30 \\ \hline
平台 & 洛谷 \\ \hline
题号 & P1005 \\ \hline
难度 & 普及 \\ \hline
标签 & DP, 高精度 \\ \hline
错题状态 & 待复习 \\ \hline
\end{tabular}

\subsection{题目描述}
对于n×m矩阵，每次从每行取首或尾元素，共取m次，每次取数乘2\^{}i，求最大总得分。

\subsection{题解}
这个题非常坑非常难，我做了好几天了。
\_\_int128
dp

\subsection{代码}
\begin{lstlisting}[language=C++]
#include <bits/stdc++.h>
using namespace std;
int main() {
    // 使用 __int128 类型
    __int128 x = 0;
    return 0;
}

\end{lstlisting}


\section{P1036【子集枚举法】}
\begin{tabular}{|l|l|}
\hline
队员 & 廖夏 \\ \hline
日期 & 2026-07-26 \\ \hline
平台 & 洛谷 \\ \hline
题号 &  \\ \hline
难度 & 普及- \\ \hline
标签 &  \\ \hline
错题状态 & 非错题 \\ \hline
\end{tabular}

\subsection{题目描述}
从n个整数中任选k个相加，求有多少种方案使得和为素数。

\subsection{题解}
核心思路：
1. 交：A1 \& A2
2. 并：A1 | A2
3. 补：A1 \^{} A2

\subsection{代码}
\begin{lstlisting}[language=C++]
#include<iostream>
int main() {
    int a = 1 & 2;
    return 0;
}

\end{lstlisting}


\section{P5143}
\begin{tabular}{|l|l|}
\hline
队员 & 廖夏 \\ \hline
日期 & 2026-07-21 \\ \hline
平台 & 洛谷 \\ \hline
题号 &  \\ \hline
难度 & 普及- \\ \hline
标签 &  \\ \hline
错题状态 & 非错题 \\ \hline
\end{tabular}

\subsection{题目描述}
爬山的题目，计算三维空间中的路径长度。

\subsection{题解}
解题思路：

\begin{lstlisting}[language=C++]
#include<iostream>
#include<algorithm>
using namespace std;
int main() { return 0; }

\end{lstlisting}

代码中注意浮点数精度。


\section{回文质数}
\begin{tabular}{|l|l|}
\hline
队员 & 廖夏 \\ \hline
日期 & 2026-07-30 \\ \hline
平台 & 洛谷 \\ \hline
题号 & P1217 \\ \hline
难度 & 普及- \\ \hline
标签 & 暴力 \\ \hline
错题状态 & 非错题 \\ \hline
\end{tabular}

\subsection{题目描述}
找出区间$[a,b]$内所有既是质数又是回文数的数字。

\subsection{题解}
力大砖飞，暴力构造。


\section{{}[NOIP 2008 提高组] 火柴棒等式}
\begin{tabular}{|l|l|}
\hline
队员 & 廖夏 \\ \hline
日期 & 2026-07-31 \\ \hline
平台 & 洛谷 \\ \hline
题号 & P1149 \\ \hline
难度 & 普及- \\ \hline
标签 & 枚举, 模拟 \\ \hline
错题状态 & 非错题 \\ \hline
\end{tabular}

\subsection{题目描述}
用n根火柴棍拼出多少个A+B=C的等式。

\subsection{题解}
注意范围，A、B不会超过1000。

\end{document}
